\documentclass[12pt,a4paper]{article}

% --- Paquetes Esenciales ---
\usepackage[utf8]{inputenc}
\usepackage[spanish, es-tabla]{babel}
\usepackage{amsmath}
\usepackage{amsfonts}
\usepackage{amssymb}
\usepackage{graphicx}
\usepackage{hyperref}
\usepackage{geometry}
\usepackage{xcolor}
\usepackage{listings}
\usepackage{tikz}
\usetikzlibrary{shapes.geometric, arrows, trees, positioning, shadows, calc}

% --- Configuración de Márgenes ---
\geometry{left=2.5cm, right=2.5cm, top=3cm, bottom=3cm}

% --- Configuración de Colores para Código ---
\definecolor{vscodeblue}{rgb}{0.0, 0.45, 0.74}
\definecolor{vscodegreen}{rgb}{0.13, 0.55, 0.13}
\definecolor{vscodepurple}{rgb}{0.58, 0.0, 0.82}
\definecolor{backcolour}{rgb}{0.98, 0.98, 0.95}

% --- Configuración de Listings para JavaScript/Bash ---
\lstdefinelanguage{JavaScript}{
  keywords={break, case, catch, continue, debugger, default, delete, do, else, false, finally, for, function, if, in, instanceof, new, null, return, switch, this, throw, true, try, typeof, var, void, while, with, let, const, async, await, yield, export, import, default, useMemo, useCallback, useEffect, useState, useContext},
  morecomment=[l]{//},
  morecomment=[s]{/*}{*/},
  morestring=[b]',
  morestring=[b]",
  ndkeywords={class, export, boolean, throw, implements, import, this},
  keywordstyle=\color{vscodeblue}\bfseries,
  ndkeywordstyle=\color{vscodepurple}\bfseries,
  identifierstyle=\color{black},
  commentstyle=\color{vscodegreen}\ttfamily,
  stringstyle=\color{orange}\ttfamily,
  sensitive=true
}

\lstset{
    language=JavaScript,
    backgroundcolor=\color{backcolour},   
    basicstyle=\ttfamily\footnotesize,
    breaklines=true,                 
    captionpos=b,                    
    numbers=left,                    
    numbersep=10pt,                  
    frame=single,
    rulecolor=\color{gray!30}
}

% --- Estilos TikZ ---
\tikzset{
    block/.style={draw, thick, text centered, minimum width=3.5cm, minimum height=1.2cm, font=\sffamily, rounded corners, drop shadow},
    provider/.style={block, fill=blue!15},
    hook/.style={block, fill=orange!15},
    ui/.style={block, fill=green!15},
    storage/.style={cylinder, draw, shape border rotate=90, minimum width=2.5cm, minimum height=2cm, fill=gray!15, font=\sffamily},
    arrow/.style={thick, ->, >=stealth}
}

\title{Especificación Técnica y Arquitectura de Software\\ \textbf{React Todo Application: Enterprise Edition}}
\author{Senior Software Architect - LJCR}
\date{15 de abril de 2026}

\begin{document}

\maketitle

\begin{abstract}
Este documento es la especificación técnica definitiva de la aplicación React Todo. Consolida el propósito del producto, la guía de despliegue, la arquitectura de componentes y el análisis de patrones de diseño. Su objetivo es servir como manual de referencia para el desarrollo, auditoría y escalado de la plataforma.
\end{abstract}

\tableofcontents
\newpage

\section{Descripción del Producto y Propósito}
La aplicación de TODOs es una solución reactiva diseñada para la gestión eficiente de tareas personales. El sistema permite la creación, filtrado, finalización y eliminación de ítems, garantizando la persistencia de datos mediante almacenamiento local asíncrono.

\section{Guía de Despliegue y Operaciones}
El proyecto utiliza \textbf{pnpm} como gestor de paquetes para optimizar el almacenamiento y la velocidad de instalación.

\begin{lstlisting}[language=bash, caption={Comandos de operación del sistema}]
# 1. Instalación de dependencias
pnpm install

# 2. Ejecución en entorno de desarrollo (Vite)
pnpm run dev

# 3. Compilación para producción
pnpm run build
\end{lstlisting}

Los archivos generados por el proceso de \texttt{build} se ubican en el directorio \texttt{dist/}, optimizados para ser servidos de forma estática.

\section{Arquitectura y Mapa del Proyecto}
La arquitectura sigue un enfoque \textbf{Feature-Based} con una clara separación entre estilos, lógica y componentes.

\subsection{Estructura de Directorios}
\begin{itemize}
    \item \texttt{src/components/}: Átomos y moléculas de la interfaz (reutilizables).
    \item \texttt{src/context/}: Capa de orquestación y estado global (Cerebro).
    \item \texttt{src/hook/}: Abstracción de lógica de persistencia (Motor).
    \item \texttt{src/style/}: Encapsulamiento de estilos CSS por componente.
\end{itemize}

\subsection{Diagrama de Jerarquía y Montado}
\begin{center}
\begin{tikzpicture}[node distance=2.2cm]
    \node (main) [block, fill=red!10] {main.jsx (EntryPoint)};
    \node (prov) [provider, below=of main] {TodoProvider (Context)};
    \node (appui) [ui, below=of prov] {AppUI (Orchestrator)};
    \node (components) [ui, below=of appui] {Componentes UI};

    \draw [arrow] (main) -- node[anchor=west, font=\tiny] {Mounts} (prov);
    \draw [arrow] (prov) -- node[anchor=west, font=\tiny] {Wraps} (appui);
    \draw [arrow] (appui) -- node[anchor=west, font=\tiny] {Renders} (components);
\end{tikzpicture}
\end{center}

\section{Gestión de Estado y Ciclo de Vida}

\subsection{Context API y Persistencia}
El uso de \textbf{TodoContext} centraliza la fuente de verdad. La persistencia se delega al hook \texttt{useLocalStorage}, el cual sincroniza el estado reactivo con el disco mediante serialización JSON.

\subsection{Secuencia de Ejecución de Tareas}
\begin{center}
\begin{tikzpicture}[node distance=3cm]
    \node (init) [circle, draw, fill=black, inner sep=2pt] {};
    \node (pend) [block, fill=yellow!10, right=of init] {Pendiente};
    \node (comp) [block, fill=blue!10, right=of pend] {Completado};
    \node (del) [block, fill=red!10, below=of comp] {Eliminado};

    \draw [arrow] (init) -- node[anchor=south, font=\tiny] {addTodo} (pend);
    \draw [arrow] (pend) -- node[anchor=south, font=\tiny] {completeTodo} (comp);
    \draw [arrow] (comp) -- node[anchor=west, font=\tiny] {deleteTodo} (del);
    \draw [arrow] (pend) |- node[anchor=north, near start, font=\tiny] {deleteTodo} (del);
\end{tikzpicture}
\end{center}

\section{Alcance y Limitaciones}
Este sistema es una demostración técnica de alta ingeniería con las siguientes restricciones actuales:
\begin{itemize}
    \item \textbf{Persistencia Local:} Datos limitados al navegador del cliente.
    \item \textbf{Single-Page:} Sin enrutamiento complejo para priorizar la velocidad de carga.
    \item \textbf{Mejoras Futuras:} Integración con Backend (API REST) y autenticación de usuarios.
\end{itemize}

\section{Conclusión}
La arquitectura robusta y el uso de estándares modernos como \textit{Vite}, \textit{Context API} y \textit{Custom Hooks} aseguran que el proyecto sea una base sólida para el escalado a nivel empresarial.

\end{document}
